\documentclass[11pt]{article}
\usepackage[T1]{fontenc}
\usepackage{lmodern,microtype,amsmath,amssymb,amsthm,mathtools}
\usepackage[margin=1in]{geometry}
\usepackage{tikz-cd,enumitem,xurl}
\usepackage[hidelinks]{hyperref}
\hypersetup{pdftitle={Application and substitution for external higher types},pdfsubject={Partial function types, computation names, and the term-substitution law}}
\setlength{\emergencystretch}{2em}
\newtheorem{theorem}{Theorem}[section]
\newtheorem{proposition}[theorem]{Proposition}
\newtheorem{corollary}[theorem]{Corollary}
\theoremstyle{definition}
\newtheorem{definition}[theorem]{Definition}
\newtheorem{example}[theorem]{Example}
\newcommand{\N}{\omega}
\newcommand{\Pcat}{\mathcal P}
\newcommand{\Vcat}{\mathcal V}
\newcommand{\dom}{\operatorname{dom}}
\newcommand{\supp}{\operatorname{supp}}
\newcommand{\id}{\operatorname{id}}
\newcommand{\ev}{\operatorname{ev}}
\newcommand{\FV}{\operatorname{FV}}
\newcommand{\pArr}{\mathbin{\Rightarrow_{\!p}}}
\newcommand{\den}[1]{\mathopen{[\mkern-3mu[}#1\mathclose{]\mkern-3mu]}}
\newcommand{\ats}{\mathord{@}}
\title{Application and substitution for external higher types}
\author{}
\date{Working note, version 1 --- 21 October 2026}
\begin{document}
\maketitle
\begin{abstract}
We extend the substitution calculus of \emph{The Unbearable Notation of
Logic} to partial computations. Strict simultaneous evaluation introduces
a domain restriction that can invalidate the ordinary substitution equation.
Occurrence maps determine this restriction and give an exact formula using
only the arguments that occur. We then construct partial function types
over partial equivalence relations on program codes. Evaluation descends
to their quotients, and the parameter theorem supplies total currying and
computable composition of function values. A monad of suspended computations
restores the term-substitution equation on computation-valued assignments.
Passing these computations as arguments gives a lazy finite-type semantics
with substitution and beta equations for arbitrary argument terms.
These constructions give a concrete semantic extension of the classical
recursion category and specify the realization problem for other recursion
categories and their coalgebraic representations.
\end{abstract}

\section{Syntax, values, and computations}

The substitution calculus in \cite[Sections 3--6 and 8--9]{Notation}
distinguishes a syntactic substitution, an assignment, and evaluation.
For a finitary signature $\Sigma$, let $T_\Sigma X$ be the set of
finite terms with variables in $X$. A substitution is a total function
$\sigma:X\to T_\Sigma Y$. Its extension
$\sigma^*:T_\Sigma X\to T_\Sigma Y$ replaces each variable $x$
by the term $\sigma(x)$, simultaneously. If $\tau:Y\to T_\Sigma Z$,
their composite is $\tau\star\sigma=\tau^*\sigma$.
Structural induction gives
$(\tau\star\sigma)^*=\tau^*\sigma^*$, the associativity law
for the term monad's Kleisli composition.

For a total $\Sigma$-algebra $A$ and an assignment
$\ats:Y\to A$, write $\ats t$ for the value of a term.
The induced assignment is
$(\ats\bullet\sigma)(x)=\ats\sigma(x)$. The substitution law is
\begin{equation}\label{eq:at-law}
 \ats(\sigma^*t)=(\ats\bullet\sigma)t.
\end{equation}
To interpret computations, one must specify what an assignment contains
and when its entries are evaluated. The distinction between values and
computations is the basis of the monadic semantics of
\cite{Moggi,BFS}. We first determine the domain correction for strict
partial evaluation, then construct function types and computation names
for which the relevant operations are explicit.

All computations below are deterministic and have nontermination as their
only effect. Composition $gf$ means first $f$, then $g$. Equality of
partial functions includes equality of their domains. For
$f:Z\rightharpoonup A$, write $\bar f:Z\rightharpoonup Z$ for
the identity on $\dom f$. A finite partial tuple
$\langle f_1,\ldots,f_k\rangle:Z\rightharpoonup A^k$ is defined
exactly where all its components are defined.

\section{The domain correction to substitution}

Let each $n$-ary symbol of $\Sigma$ denote a partial function
$f^A:A^n\rightharpoonup A$. Evaluation is strict: it evaluates all
immediate subterms before applying $f^A$. This gives each term $t$ over
a finite context $X$ a partial interpretation
$\den t_X:A^X\rightharpoonup A$. Nullary symbols may also diverge.

\begin{example}\label{ex:weakening}
Take $A=\N$, a constant $0$, a divergent constant $\Omega$, and
a binary symbol $p$ whose value operation is $p(a,b)=a$. On total
assignments, the terms $x$ and $p(x,y)$ denote the same function
$A^2\to A$. Under the substitution
$\sigma(x)=0$, $\sigma(y)=\Omega$, the first term returns $0$
and the second diverges. Thus equality on total value assignments is
not a congruence for arbitrary computational substitution.
\end{example}

The same issue occurs before any quotient is formed. Evaluating every
entry of the substituted assignment produces the nowhere-defined tuple
$\langle0,\Omega\rangle$, even when the outer term is just $x$.

\begin{proposition}\label{prop:restriction-substitution}
Let $t\in T_\Sigma X$ and $\sigma:X\to T_\Sigma Y$, where
$X,Y$ are finite. Put
\[
 s=\langle\den{\sigma(x)}_Y\rangle_{x\in X},
 \qquad g=\den{\sigma^*t}_Y.
\]
Then
\begin{equation}\label{eq:restriction-substitution}
 \den t_X\,s=g\bar s.
\end{equation}
If $S=\FV(t)$ and
$s_S=\langle\den{\sigma(x)}_Y\rangle_{x\in S}$, then
\begin{equation}\label{eq:support-substitution}
 \den t_S\,s_S=\den{\sigma^*t}_Y.
\end{equation}
\end{proposition}
\begin{proof}
On $\dom s$, every substituted term has a value. Structural induction
on $t$ proves that evaluation after substitution and evaluation under
this value assignment have the same definedness and value. Outside
$\dom s$, both sides of \eqref{eq:restriction-substitution} are
undefined. This proves the first identity.

For the second, apply the first identity in context $S$. If
$\sigma^*t$ converges, every leaf substitution corresponding to an
occurrence of a variable in $t$ converges, by strictness. Hence
$\dom g\subseteq\dom s_S$, and $g\bar s_S=g$. The empty
support case uses the unique total map into the singleton $A^\varnothing$.
\end{proof}

The categorical calculation uses
$\pi_xs=\den{\sigma(x)}\bar s$ and
$\langle g_i e\rangle=\langle g_i\rangle e$ for a nonempty tuple
and a restriction idempotent $e$. For the nullary step, the restriction
terminal maps satisfy $!_X s=!_Y\bar s$.
The same structural induction proves the proposition
in a cartesian restriction category, using restriction products and
a restriction terminal object for the empty context.

\subsection{Occurrence maps and evaluation order}

The occurrence construction of \cite[Section 9 and Appendix A.5]{Notation}
records a map $o:\{1,\ldots,m\}\to\{1,\ldots,k\}$: the
$r$th occurrence uses argument $o(r)$. It induces the total routing map
$o^*:A^k\to A^m$. For partial maps $f_i:Z\rightharpoonup A$,
\begin{equation}\label{eq:routing}
 o^*\langle f_1,\ldots,f_k\rangle
 =\langle f_{o(1)},\ldots,f_{o(m)}\rangle
    \overline{\langle f_1,\ldots,f_k\rangle}.
\end{equation}
Both sides evaluate the selected coordinates on the intersection of
all $k$ component domains. This identity explains why routing after
strict tuple evaluation cannot discard an unused divergence.

Let $t_{\rm lin}$ be obtained by giving each variable occurrence of
$t$ a fresh variable. With $x_1,\ldots,x_k$ the declared variables,
the exact occurrence formula is
\begin{equation}\label{eq:linear-substitution}
 \den{t[\sigma]}
 =\den{t_{\rm lin}}
   \langle\den{\sigma(x_{o(1)})},\ldots,
           \den{\sigma(x_{o(m)})}\rangle.
\end{equation}
Structural induction proves the equality, including domains. Repeated
occurrences evaluate the same deterministic computation and obtain the
same value. The fibrewise occurrence maps of \cite{Notation} record
the routing in iterated substitution. These formulas concern algebraic
terms with strict operations; abstraction will require the distinction
between a function value and a computation producing one.

\section{Partial function types on program quotients}\label{sec:per}

Fix an acceptable numbering $(\varphi_e)_{e\in\N}$ of the unary
partial recursive functions. Use the Cantor pairing bijection
\[
 \langle m,n\rangle=\frac{(m+n)(m+n+1)}2+n
 \quad(m,n\in\N),
\]
whose two projections are computable.
The parameter theorem supplies a total recursive function $s$ such that
\begin{equation}\label{eq:smn}
 \varphi_{s(e,t)}(a)\simeq\varphi_e(\langle t,a\rangle).
\end{equation}
Here $\simeq$ means that both sides diverge or both converge to the
same number. The code $s(e,t)$ is produced without running the
computation on the right.

A partial equivalence relation, or PER, is a symmetric transitive
relation $R\subseteq\N^2$. Its support is
$\supp R=\{a:aRa\}$, and its represented set is
$|R|=\supp R/R$. We allow arbitrary PERs; their supports and
relations need not be computably enumerable. Let $\mathbf N$ be
equality on $\N$. Products are coded by pairing:
$\langle a,b\rangle\,(R\times S)\,\langle a',b'\rangle$
exactly when $aRa'$ and $bSb'$.

\begin{definition}
An arrow $f:R\rightharpoonup S$ is a partial function
$|R|\rightharpoonup|S|$ with a program $e$ such that, for every
$a\in\supp R$,
\[
 f([a]_R)\downarrow\quad\Longleftrightarrow\quad
       \varphi_e(a)\downarrow,
\]
and, when defined, $\varphi_e(a)\in\supp S$ and
$f([a]_R)=[\varphi_e(a)]_S$. Such a program \emph{tracks} $f$.
Write $\Pcat$ for the category of PERs and these arrows, and
$\Vcat=\operatorname{Tot}(\Pcat)$ for its total subcategory.
\end{definition}

The definition requires definedness to be invariant under the input PER.
Identity is tracked by the identity program, and program composition
tracks composition of arrows. Equality of arrows is equality of the
represented partial functions, so composition is associative.
Restriction is tracked by running the tracker and returning the original
input when it halts. Pairing and projections give restriction products.
A PER with support $\{0\}$ gives a restriction terminal object.
Thus $\Pcat$ is a cartesian restriction category. Its full subcategory
on the positive finite powers of $\mathbf N$ is the classical
recursion category, up to the chosen numerical codings.

\begin{definition}\label{def:partial-function-per}
For PERs $R,S$, define $R\pArr S$ by declaring
$e\,(R\pArr S)\,e'$ if, for every $aRa'$, either
\begin{enumerate}[nosep,label=(\roman*)]
\item both $\varphi_e(a)$ and $\varphi_{e'}(a')$ diverge; or
\item both converge and their values are $S$-related.
\end{enumerate}
\end{definition}

Symmetry follows by interchanging $a,a'$; transitivity follows by using
the middle program at $a'$ and the fact that $a'Ra'$. Thus this is a
PER. Its support consists exactly of the trackers of arrows
$R\rightharpoonup S$, and its quotient identifies precisely those
trackers that induce the same arrow. In particular,
$\supp(\mathbf N\pArr\mathbf N)=\N$: an arbitrary numerical
program denotes a partial numerical function.

\begin{theorem}\label{thm:partial-currying}
Universal application induces an arrow
\[
 \ev_{R,S}:(R\pArr S)\times R\rightharpoonup S,
 \qquad \ev([e],[a])=[\varphi_e(a)]
\]
on its domain of convergence. For every PER $Q$, there is a bijection
\begin{equation}\label{eq:currying}
 \Pcat(Q\times R,S)\ \cong\ \Vcat(Q,R\pArr S),
 \qquad h\longmapsto\Lambda h,
\end{equation}
natural under total changes of parameter $Q$. Its equations are
\begin{equation}\label{eq:beta-eta}
 \ev(\Lambda h\times\id_R)=h,
 \qquad \Lambda\bigl(\ev(g\times\id_R)\bigr)=g
       \quad(g\text{ total}).
\end{equation}
\end{theorem}
\begin{proof}
Definition~\ref{def:partial-function-per} makes the evaluator's
domain saturated under $(R\pArr S)\times R$ and its output class
independent of representatives. The universal program tracks it.

If $e$ tracks $h$, define $\Lambda h$ on names by
$t\mapsto s(e,t)$. For $tQt'$ and $aRa'$, the tracking condition
for $h$ gives matching definedness and $S$-related outputs on
$\langle t,a\rangle$ and $\langle t',a'\rangle$.
Equation~\eqref{eq:smn} therefore gives
$s(e,t)\,(R\pArr S)\,s(e,t')$. The compiler is total, so it
tracks a total arrow $Q\to R\pArr S$.

The first equation in \eqref{eq:beta-eta} follows from
\eqref{eq:smn}. Conversely, if $g$ is total, its value at $[t]$
is a class of programs. Uncurrying and currying return a program
with the same partial behavior on every valid $R$-input, and hence
the same class. This proves the second equation and uniqueness.
For total $u:Q'\to Q$, both $\Lambda(h(u\times\id_R))$
and $(\Lambda h)u$ classify the same family of partial functions,
proving naturality.
\end{proof}

\subsection{Composition and the partial-parameter boundary}

A total recursive compiler $c$ for composition satisfies
\[
 \varphi_{c(d,e)}(a)\simeq\varphi_d(\varphi_e(a)),
\]
where the right side first runs $e$. It induces a total arrow
\[
 (S\pArr U)\times(R\pArr S)\longrightarrow R\pArr U.
\]
An identity-program code gives the identity function value. These
operations satisfy the unit and associativity equations on the
quotients. Program indices for different compilation orders need not
be equal.

The restriction to total parameters in Theorem~\ref{thm:partial-currying}
is exact. Suppose $|Q'|$ is nonempty and
$u:Q'\rightharpoonup Q$ is nowhere defined. Then
$(\Lambda h)u$ is nowhere defined. But
$\Lambda(h(u\times\id_R))$ is total and returns the class of the
everywhere-divergent function. A divergent computation of a function
value and a function value that diverges on every argument are distinct.
The next construction preserves this distinction during substitution.

\section{Computation names and the substitution law}\label{sec:lift}

For a PER $S$, define a lifting PER $LS$ on program indices by
\[
 e\,(LS)\,e'\quad\Longleftrightarrow\quad
 \begin{cases}
 \varphi_e(0),\varphi_{e'}(0)\text{ both diverge};\quad\text{or}\\
 \varphi_e(0),\varphi_{e'}(0)\text{ both converge to $S$-related values}.
 \end{cases}
\]
A valid name is a suspended computation that either diverges or returns
a valid $S$-name. The quotient $|LS|$ consists of a bottom element and
one element for each member of $|S|$. Naming a computation does not
decide whether it will converge.

\begin{proposition}\label{prop:lift}
The construction $L$ is a monad on $\Vcat$, and there are natural
bijections
\begin{equation}\label{eq:kleisli-partial}
 \Pcat(R,S)\cong\Vcat(R,LS)
\end{equation}
identifying $\Pcat$ with its Kleisli category.
\end{proposition}
\begin{proof}
For $f:R\rightharpoonup S$ tracked by $e$, the parameter theorem
produces, uniformly and totally in $a$, a code for the suspended
computation $\varphi_e(a)$. This defines a total arrow
$\widehat f:R\to LS$. Conversely, a total arrow $R\to LS$
can be followed by forcing the returned computation name on input
$0$. The definition of $LS$ supplies precisely the invariant
definedness and values required of an arrow in $\Pcat$. These
constructions are inverse on represented functions.

For total $f:S\to T$, $Lf$ sends a computation name to a new name
that runs it and then applies a tracker of $f$. The unit
$\eta_S:S\to LS$ returns a code for the computation immediately
returning its input name. The multiplication $\mu_S:LLS\to LS$
returns a code that runs the outer computation and then the resulting
inner computation. All these code-producing programs are total on
valid inputs, by the parameter theorem. Independence of trackers
follows from equality of the represented results and their definedness.

The monad laws follow by comparing executions: returning before or
after a computation preserves its convergence and value; either
parenthesization of flattening runs the same three successive
computations. Kleisli composition runs the first computation and then
the second, so \eqref{eq:kleisli-partial} preserves identities and
composition.
\end{proof}

The monad has computable strength
$R\times LS\to L(R\times S)$: retain the $R$-name and return
a suspended computation that forces the $S$-name and pairs the result
with it. This gives the value/computation semantics of
\cite{Moggi,BFS} concretely on the program quotients.
Combining \eqref{eq:currying} and \eqref{eq:kleisli-partial} gives
\[
 \Vcat(Q\times R,LS)\cong\Vcat(Q,R\pArr S).
\]
Thus $R\pArr S$ represents the function space from values of type
$R$ to computations of type $S$.

\subsection{Restoring the paper's equation}

Let $A$ be a PER and interpret each symbol $f$ of $\Sigma$ by an
arrow $f:A^n\rightharpoonup A$ of $\Pcat$. Define a total
operation $f^L:(LA)^n\to LA$ as follows: from $n$ computation
names, construct a suspended computation that forces them in order
and then runs a tracker of $f$ on the resulting values. The program
constructing this name is total. The represented computation may
diverge. For a nullary operation it simply suspends that operation.

These operations give the underlying set $|LA|$ an ordinary total
$\Sigma$-algebra structure. Assignments now have type
$\ats:X\to|LA|$: every variable is assigned a computation value,
including possibly bottom.

\begin{theorem}\label{thm:computation-substitution}
For this algebra and every syntactic substitution
$\sigma:X\to T_\Sigma Y$,
\[
 \ats(\sigma^*t)=(\ats\bullet\sigma)t
 \qquad(t\in T_\Sigma X).
\]
Equality is equality of computations: both diverge, or both return
the same $A$-class. If the signature and its primitive trackers are
effective, the construction is uniformly computable from a finite term,
substitution terms for its variables, and computation names for the
finitely many assignment entries used.
\end{theorem}
\begin{proof}
The variable case is the definition of $\ats\bullet\sigma$.
At a node $f(t_1,\ldots,t_n)$, apply the induction hypotheses to
each subterm and then the total operation $f^L$. This proves the
equation by structural induction, exactly as for the total term
algebra in \cite{Notation}. Recursive syntax processing and the
total suspension compilers prove the effective assertion.
\end{proof}

An unused divergent computation is now an available assignment entry
whose name is never forced. The distinction underlying
Example~\ref{ex:weakening} is
\[
 (|A|_\bot)^X\quad\text{versus}\quad(|A|^X)_\bot.
\]
The first set records partial information coordinate by coordinate;
the second records only success or failure of the whole tuple. The
algebra on $|LA|$ distinguishes $x$ from strict $p(x,y)$ and makes
equality on all computation-valued assignments a substitution
congruence. The equality is relative to this interpretation of the
primitive symbols. In particular, strict lifting of the numerical
projection $p(a,b)=a$ forces both input computations. It is different
from the structural projection $(LA)^2\to LA$, which simply selects
the first computation value. The equation $p(x,y)=x$ of the original
value algebra therefore need not hold in the lifted algebra.

\section{Higher-type application and strict sequencing}

Starting with $A_\iota=\mathbf N$, define finite value types by
\[
 A_{\sigma\times\tau}=A_\sigma\times A_\tau,
 \qquad A_{\sigma\to\tau}=A_\sigma\pArr A_\tau.
\]
Computations of type $\sigma$ are represented by $LA_\sigma$.
Application is the partial arrow
$A_{\sigma\to\tau}\times A_\sigma\rightharpoonup A_\tau$
of Theorem~\ref{thm:partial-currying}. Suspending it gives the
total value-category arrow
\[
 A_{\sigma\to\tau}\times A_\sigma\longrightarrow LA_\tau.
\]
A body $h:\Gamma\times A_\sigma\rightharpoonup A_\tau$
has a total abstraction $\Lambda h:\Gamma\to A_{\sigma\to\tau}$.
These provide partial application and abstraction at every finite type.

Substitution that first computes a value context is strict sequencing. If
$u:\Delta\rightharpoonup\Gamma$ and
$t:\Gamma\rightharpoonup A$, then the computation
\emph{let $\gamma\leftarrow u$ in $t$} is $tu$. Its domain is
\[
 \{d\in\dom u:u(d)\in\dom t\}.
\]
On suspended computations this is monadic bind. For a parameterized
body $h:\Gamma\times R\rightharpoonup S$ and an argument
$u:\Gamma\rightharpoonup R$, the exact application equation is
\begin{equation}\label{eq:cbv-beta}
 \ev\langle\Lambda h,u\rangle
       =h\langle\id_\Gamma,u\rangle.
\end{equation}
Both sides first require $u$ to converge, and then apply the same
partial function to its value. If $h$ ignores its second argument
but $u$ diverges, both sides diverge. For a total change of context
$v:\Delta\to\Gamma$, Theorem~\ref{thm:partial-currying} gives
\[
 (\Lambda h)v=\Lambda\bigl(h(v\times\id_R)\bigr).
\]
Thus total value substitution commutes with abstraction, while
substitution that first executes a computation is represented by
explicit sequencing.

This gives two distinct associative compositions. Syntactic
substitution is Kleisli composition for the term monad $T_\Sigma$.
Execution and strict sequencing are Kleisli composition for the
partiality monad $L$. Theorem~\ref{thm:computation-substitution}
connects term substitution to evaluation on computation names;
equation~\eqref{eq:cbv-beta} describes application after a computation
produces an argument value. Their types determine which law applies.

\section{Lazy application and substitution}\label{sec:lazy}

Lazy application passes a computation name to the function without
first forcing it. For PERs $A,B$, its function values therefore have
type $(LA)\pArr B$, rather than $A\pArr B$. The evaluator
\[
 \ev:((LA)\pArr B)\times LA\rightharpoonup B
\]
may return a value even when its second argument represents a divergent
computation. A program returning $0$ without inspecting its input
is an example when $B=\mathbf N$.

For a body $h:\Gamma\times LA\rightharpoonup B$ and an argument
$u:\Gamma\rightharpoonup A$, suspension gives a total arrow
$\widehat u:\Gamma\to LA$. The application law is
\begin{equation}\label{eq:lazy-body-beta}
 \ev\langle\Lambda h,\widehat u\rangle
       =h\langle\id_\Gamma,\widehat u\rangle.
\end{equation}
Unlike \eqref{eq:cbv-beta}, this equation does not first execute $u$.
The function body decides whether to force its computation name.

\subsection{A finite-type semantics}

Consider types generated by a numerical base type $\iota$ and arrows.
Define value PERs $V_\sigma$ and computation PERs $C_\sigma$ by
\[
 V_\iota=\mathbf N,\qquad C_\sigma=LV_\sigma,\qquad
 V_{\sigma\to\tau}=C_\sigma\pArr V_\tau.
\]
For a finite context
$\Gamma=(x_1:\sigma_1,\ldots,x_k:\sigma_k)$, put
$E_\Gamma=\prod_{i=1}^k C_{\sigma_i}$, with the terminal PER
for the empty context. An environment contains computation values;
its formation does not execute the computations they name.

There is a total arrow
\begin{equation}\label{eq:lazy-app}
 a_{\sigma,\tau}:C_{\sigma\to\tau}\times C_\sigma\to C_\tau.
\end{equation}
On input names $p,q$, its tracker returns a suspended program which
runs $\varphi_p(0)$ to obtain a function index $e$, and then runs
$\varphi_e(q)$. It passes $q$ itself to $e$, without first running
$\varphi_q(0)$. The parameter theorem makes the code-producing map
total. The definitions of the lifting and function PERs prove
independence of representatives. Application forces the function
computation; the function controls evaluation of its argument.

Use the simply typed lambda calculus with variables, abstraction,
application, numerical constants, and a divergent constant
$\Omega_\sigma$ at each type. Interpret a judgement
$\Gamma\vdash M:\tau$ by a total arrow
\[
 \den M_\Gamma:E_\Gamma\to C_\tau
\]
of $\Vcat$. Variables are projections. A numeral denotes the
constant immediately returning that number, and $\Omega_\sigma$
denotes the constant bottom computation. Application is interpreted by
\[
 \den{MN}_\Gamma
 =a_{\sigma,\tau}\langle\den M_\Gamma,\den N_\Gamma\rangle.
\]
For a body $\Gamma,x:\sigma\vdash M:\tau$, let
$h:E_\Gamma\times C_\sigma\rightharpoonup V_\tau$ be the
partial arrow corresponding to $\den M_{\Gamma,x}$ under
\eqref{eq:kleisli-partial}. Define
\begin{equation}\label{eq:lazy-lambda}
 \den{\lambda x.M}_\Gamma
       =\eta_{V_{\sigma\to\tau}}\Lambda h.
\end{equation}
Thus abstraction immediately returns a function value whose body can
later force an argument computation. These clauses construct a tracker
uniformly from a finite typed term.

\begin{theorem}[Lazy substitution and beta]\label{thm:lazy-substitution}
Let $\theta$ assign to every $x_i:\sigma_i$ in $\Gamma$ a term
$\Delta\vdash\theta(x_i):\sigma_i$. Let $M[\theta]$ denote
simultaneous capture-avoiding substitution, and put
\[
 s_\theta=\langle\den{\theta(x_1)}_\Delta,\ldots,
                       \den{\theta(x_k)}_\Delta\rangle:
 E_\Delta\to E_\Gamma.
\]
For $\Gamma\vdash M:\tau$,
\begin{equation}\label{eq:lazy-substitution}
 \den{M[\theta]}_\Delta=\den M_\Gamma s_\theta.
\end{equation}
For $\Gamma,x:\sigma\vdash M:\tau$ and
$\Gamma\vdash N:\sigma$,
\begin{equation}\label{eq:lazy-beta}
 \den{(\lambda x.M)N}_\Gamma
       =\den{M[N/x]}_\Gamma.
\end{equation}
No convergence assumption on $N$ is needed.
\end{theorem}
\begin{proof}
Every component of $s_\theta$ is total: it returns a computation name,
including a name for bottom. The variable case of
\eqref{eq:lazy-substitution} is therefore the projection identity.
The constant cases hold because reindexing is total. The application
case follows by induction, pairing, and composition with
$a_{\sigma,\tau}$.

For abstraction, rename the bound variable to avoid capture. Extension
of $\theta$ by that variable induces the total map
$s_\theta\times\id_{C_\sigma}$. By induction and the naturality
of \eqref{eq:kleisli-partial}, the substituted body corresponds to
$h(s_\theta\times\id_{C_\sigma})$. Total-parameter naturality
in Theorem~\ref{thm:partial-currying} and
\eqref{eq:lazy-lambda} give the required abstraction equation.

For beta, put $q=\langle\id_{E_\Gamma},\den N_\Gamma\rangle$.
Forcing the immediately returned function in
\eqref{eq:lazy-lambda} and applying it to the argument name gives
\[
 a_{\sigma,\tau}
   \langle\eta_{V_{\sigma\to\tau}}\Lambda h,\den N_\Gamma\rangle
 =\widehat{hq}
 =\den M_{\Gamma,x}q.
\]
The substitution equation identifies the last expression with
$\den{M[N/x]}_\Gamma$.
\end{proof}

In particular, $(\lambda x.0)\Omega_\sigma$ denotes the computation
returning $0$. A divergent function computation and an immediately
returned divergent function remain different: $\Omega_{\sigma\to\tau}$
diverges, whereas $\lambda x.\Omega_{\sigma\to\tau}x$ returns a
function value. The eta equation in
Theorem~\ref{thm:partial-currying} applies to total families of function
values, as its statement requires.

The same principle extends the algebraic substitution calculus in
Section~\ref{sec:lift}. One may interpret a primitive symbol by any
specified total arrow $(LA)^n\to LA$, including a non-strict one.
The proof of Theorem~\ref{thm:computation-substitution} still applies.
Such an arrow $F$ extends a given value operation
$f:A^n\rightharpoonup A$ when
$F\eta_A^n=\widehat f:A^n\to LA$.
For example, the structural projection selects its first computation
without forcing the second and satisfies $p(x,y)=x$ on all assignments.
A value operation $A^n\rightharpoonup A$ alone does not determine
this behavior on divergent arguments; the operation on computation
values specifies it.

\subsection{Sharing and occurrence maps}

Passing a suspended argument is call-by-name. Call-by-need additionally
shares its evaluation: variable occurrences referring to the same
binding use a single stored computation, whose result is retained once
obtained. For the pure lambda calculus, call-by-name and call-by-need
induce the same observational equivalence
\cite[Theorem 32]{MOW}. The quotient semantics above records
convergence and results. To describe sharing itself, an operational
presentation must also record the shared bindings.

The occurrence maps of \cite[Section 9]{Notation} specify which uses
belong to the same binding. An occurrence map
$o:\{1,\ldots,m\}\to\{1,\ldots,k\}$ can route each occurrence
to one of $k$ suspended bindings. An empty fibre causes no demand;
several occurrences in one fibre can share a result. The maps supply
the incidence data. A call-by-need evaluator must in addition implement
the store and its update rule. This distinguishes the substitution
theorem from an assertion about the cost of executing substitutions.

\section{The remaining representation problem}

The construction above provides a semantic category containing the
classical recursion category, together with explicit application,
abstraction, composition compilers, and substitution laws. The use of
arbitrary PERs supplies quotient types beyond the computably enumerable
supports considered in the regular exact completion of
\cite[section ``Total function spaces and finite observations'']{External}.
For example, the partial numerical function type has all indices as
support, while its subspace of total functions has support
$\mathrm{TOT}=\{e:\varphi_e\text{ is total}\}$.
The constructions belong to the realizability approach to higher-type
computability surveyed in \cite{Longley}.

\begin{proposition}\label{prop:no-least-range}
There is an arrow $f$ of $\Pcat$ for which no least restriction
idempotent $r$ satisfies $rf=f$.
\end{proposition}
\begin{proof}
Choose a non-computably-enumerable set $D\subseteq\N$ containing
$0$. Let $S$ have full support, with $D$ as one equivalence class
and every $n\notin D$ as a singleton class. The total arrow
$f:1\to S$ selecting $[0]_S$ is tracked by the constant program.
For each $n\notin D$, the partial identity $r_n:S\rightharpoonup S$
undefined only at $[n]_S$ is tracked by a program that diverges on
input $n$ and otherwise returns its input. Thus $r_nf=f$.

If a least $r$ existed, it would satisfy $r\leq r_n$ for every
$n\notin D$, where restriction idempotents are ordered by inclusion
of domains. Its domain would therefore consist exactly of $[0]_S$.
Since $S$ has full numerical support, any tracker of $r$ would halt
exactly on $D$, a contradiction.
\end{proof}

Consequently, $\Pcat$ does not admit ranges for all arrows, when ranges
are least restriction idempotents supporting an arrow. A construction
retaining that structure must choose a smaller family of quotients
or a different semantic extension.

For an arbitrary recursion category, the next theorem must provide
comparable computation names and total parameterization compilers in
a specified semantic extension. A coalgebraic realization must in
addition interpret evaluation and these compilers and prove compatibility
with the quotient relations and their domains. In particular, the
application square must express equality of partial maps, including
definedness; agreement only where both sides happen to converge is
insufficient. The domain and value conditions are precisely those in
the partial-application descent proposition of
\cite[section ``Application on quotient types'']{External}.

The classical model supplies concrete operations against which such
a realization can be tested. The required assertions are preservation
of application \eqref{eq:beta-eta}, total parameter substitution,
strict sequencing \eqref{eq:cbv-beta}, lazy application and
substitution in Theorem~\ref{thm:lazy-substitution}, and the
computation-valued term-substitution equation. Establishing these for a family of
coalgebraic representations would extend the classification of their
underlying recursion categories to a comparison of their higher-type
operations.

\begin{thebibliography}{9}
\small
\bibitem{Notation} F.~Lengyel, \emph{The Unbearable Notation of Logic},
version 11, 2026.
\url{https://github.com/flengyel/The_Unbearable_Notation_Of_Logic}.
\bibitem{External} F.~Lengyel, \emph{An External Finite-Type Completion
of Recursion Categories}, version 0.11, 2026.
\url{https://github.com/flengyel/External_Finite_Type_Completion}.
\bibitem{Moggi} E.~Moggi, \emph{Notions of computation and monads},
Information and Computation \textbf{93} (1991), 55--92.
\href{https://doi.org/10.1016/0890-5401(91)90052-4}{doi:10.1016/0890-5401(91)90052-4}.
\bibitem{BFS} A.~Bucalo, C.~F\"uhrmann, and A.~Simpson,
\emph{An equational notion of lifting monad},
Theoretical Computer Science \textbf{294} (2003), 31--60.
\href{https://doi.org/10.1016/S0304-3975(01)00243-2}{doi:10.1016/S0304-3975(01)00243-2}.
\bibitem{Longley} J.~R. Longley, \emph{Notions of computability at higher
types I}, in \emph{Logic Colloquium 2000}, 2005, pp.~32--142.
\href{https://doi.org/10.1017/9781316755884.004}{doi:10.1017/9781316755884.004}.
\bibitem{MOW} J.~Maraist, M.~Odersky, and P.~Wadler,
\emph{The call-by-need lambda calculus}, Journal of Functional
Programming \textbf{8} (1998), 275--317.
\href{https://doi.org/10.1017/S0956796898003037}{doi:10.1017/S0956796898003037}.
\end{thebibliography}
\end{document}
